Vier Kugeln A bis D rollen nacheinander über die Kante eines \hTO\ hohen Tischs und fliegen dann bis zum Boden. Ihre Massen und ihre Geschwindigkeiten an der Tischkante (horizontal):
\begin{center}
\begin{tabular}{lcccc}
 Kugel & A & B & C & D \\ \hline
 Masse & \MaO & \MbO & \McO & \MdO \\
 Geschwindigkeit & \vaO & \vbO & \vcO & \vdO
\end{tabular}
\end{center}
Der Luftwiderstand ist vernachlässigbar.

\begin{abcliste}
\abc
Ordne die Kugeln nach ihrer Flugzeit, von der längsten zur kürzesten. Gleiche Zeiten bekommen denselben Rang. Begründe.

\abc
Ordne die Kugeln danach, wie weit vom Tisch entfernt sie landen, von der weitesten zur nächsten. Begründe.

\abc
Berechne die Flugzeit und die Landeweiten. Rechne mit \gO.
\end{abcliste}
